\documentclass[11pt,letterpaper]{article}
\usepackage[margin=2.5cm]{geometry}
\usepackage[utf8]{inputenc}
\usepackage[T1]{fontenc}
\usepackage{lmodern}
\usepackage{amsmath,amssymb,amsthm,mathtools,bm}
\usepackage{enumitem}
\usepackage{booktabs}
\usepackage{tikz}
\usetikzlibrary{arrows.meta,patterns,positioning,calc}
\usepackage{pgfplots}
\pgfplotsset{compat=1.16}
\usepackage{titlesec}
\usepackage[hidelinks]{hyperref}

\setlength{\parindent}{0pt}
\setlength{\parskip}{0.5em}
\titleformat{\section}{\large\bfseries}{\thesection.}{0.6em}{}
\titleformat{\subsection}{\normalsize\bfseries}{\thesubsection.}{0.6em}{}

\theoremstyle{definition}
\newtheorem{definition}{Definition}[section]
\newtheorem{example}{Example}[section]
\newtheorem{exercise}{Exercise}[section]
\newtheorem{theorem}{Theorem}[section]
\newtheorem{remark}{Remark}[section]

\newcommand{\R}{\mathbb{R}}
\newcommand{\N}{\mathbb{N}}
\newcommand{\Z}{\mathbb{Z}}
\newcommand{\U}{\mathcal{U}}
\newcommand{\E}{\mathbb{E}}
\newcommand{\Var}{\operatorname{Var}}
\newcommand{\Cov}{\operatorname{Cov}}
\newcommand{\Normal}{\mathcal{N}}
\newcommand{\dd}{\,\mathrm{d}}
\newcommand{\T}{^{\top}}
\renewcommand{\P}{P}
\newcommand{\set}[1]{\left\{#1\right\}}
\newcommand{\given}{\mid}
\newcommand{\note}[1]{\textit{#1}}

\title{Lecture 9: Variational Inference}
\author{}\date{}

\begin{document}
\maketitle

%=====================================================
\section{The problem}
As always, we want to solve
\[
p(\theta \given y) = \frac{p(y \given \theta)\,p(\theta)}{p(y)},
\qquad
p(y) = \int p(y \given \theta')\,p(\theta')\dd\theta',
\]
where $p(\theta \given y) = p(\theta_1, \dots, \theta_K \given y_1, \dots, y_N)$.

The main idea behind \emph{variational inference} is to use a family of distributions over $\theta$,
\[
q_\phi(\theta),
\]
with parameters $\phi$. We need to find the $\phi$ that makes $q_\phi(\theta)$ the closest to $p(\theta \given y)$.

%=====================================================
\section{Reverse and forward KL divergence}
The true posterior is typically not in the family. A way of measuring how close $q_\phi(\theta)$ is to $p(\theta \given y)$ is the \emph{reverse} Kullback--Leibler (KL) divergence
\[
D_{\mathrm{KL}}\bigl(q_\phi(\theta) \,\|\, p(\theta \given y)\bigr) = \int q_\phi(\theta)\log\left(\frac{q_\phi(\theta)}{p(\theta \given y)}\right)\dd\theta.
\]
Since $q_\phi(\theta)$ is a pdf (or pmf), we have
\[
D_{\mathrm{KL}}(q \,\|\, p) = \E_{\theta \sim q_\phi(\theta)}\left[\log\frac{q_\phi(\theta)}{p(\theta \given y)}\right].
\]

\begin{remark}
$D_{\mathrm{KL}}$ is not an actual metric: it is not symmetric and does not satisfy the triangle inequality.
\end{remark}

Notice that the \emph{direct} (forward) KL divergence is
\[
D_{\mathrm{KL}}(p \,\|\, q) = \E_{\theta \sim p(\theta \given y)}\left[\log\frac{p(\theta \given y)}{q_\phi(\theta)}\right],
\]
which implies sampling $\theta$ from $p(\theta \given y)$, which is intractable.

%=====================================================
\section{The evidence lower bound (ELBO)}
We actually cannot minimize $D_{\mathrm{KL}}(q \,\|\, p)$ directly either. Notice that, using
\[
p(\theta \given y) = \frac{p(\theta, y)}{p(y)},
\]
we have
\[
\log\left(\frac{q_\phi(\theta)\,p(y)}{p(\theta, y)}\right) = \log q_\phi(\theta) - \log p(\theta, y) + \log p(y),
\]
and therefore
\begin{align*}
D_{\mathrm{KL}}(q \,\|\, p) &= \E_q[\log q_\phi(\theta)] - \E_q[\log p(\theta, y)] + \E_q[\log p(y)] \\
&= -\Bigl(\E_q[\log p(\theta, y)] - \E_q[\log q_\phi(\theta)]\Bigr) + \log p(y).
\end{align*}
We look for
\[
q^*_\phi(\theta) = \operatorname*{arg\,min}_\phi D_{\mathrm{KL}}(q \,\|\, p).
\]
The last term does not depend on $\phi$ (and $\E_q[\log p(y)] = \log p(y)$), so instead of minimizing the KL divergence we can maximize the quantity in the parentheses.

Notice that
\[
p(y) = \int p(\theta, y)\dd\theta,
\]
and
\[
p(y) = \int p(\theta, y)\,\frac{q_\phi(\theta)}{q_\phi(\theta)}\dd\theta = \E_q\left[\frac{p(\theta, y)}{q_\phi(\theta)}\right],
\]
so
\[
\log p(y) = \log\left(\E_q\left[\frac{p(\theta, y)}{q_\phi(\theta)}\right]\right).
\]
By Jensen's inequality for a concave function $f$ (here $f = \log$),
\[
f(\E[X]) \ge \E[f(X)],
\]
we have that
\begin{align*}
\log p(y) &\ge \E_q\left[\log\left(\frac{p(\theta, y)}{q_\phi(\theta)}\right)\right] \\
&= \E_q[\log p(\theta, y)] - \E_q[\log q_\phi(\theta)],
\end{align*}
that is,
\[
\boxed{\log p(y) \ge \mathrm{ELBO}(q)},
\]
the \emph{Evidence Lower Bound} (ELBO).

%=====================================================
\section{Mean-field variational inference}
In the mean-field variational family the variables are mutually independent:
\[
q_\phi(\theta) = q(\theta_1, \dots, \theta_K) = \prod_{i=1}^{K} q_{\phi_i}(\theta_i).
\]
This is not a model of the data. So, we want $\phi$ such that
\[
q^*_\phi(\theta) = \operatorname*{arg\,max}_\phi\, \mathrm{ELBO}(q).
\]
Several algorithms are used for maximizing the ELBO:
\begin{itemize}
  \item Coordinate ascent.
  \item Stochastic variational inference.
  \item Automatic differentiation.
\end{itemize}

%=====================================================
\section{ADVI}
Let us choose a mean-field Gaussian family:
\[
q_\phi(\theta) = \prod_{j=1}^{K}\Normal(\theta_j \given \mu_j, \sigma_j^2),
\]
where $\phi = (\mu_{1:K}, \sigma_{1:K})$. So, we need to maximize the ELBO,
\[
\mathcal{L}(\phi) = \E_{q_\phi}\bigl[\log p(\theta, y) - \log q_\phi(\theta)\bigr].
\]
We need to find
\[
\phi^* = \operatorname*{arg\,max}_\phi\, \mathcal{L}(\phi).
\]

\subsection{Score-function gradient of the ELBO}
Thus, we need to calculate gradients:
\[
\nabla_\phi\mathcal{L}(\phi) = \nabla_\phi\int q_\phi(\theta)\bigl[\log p(\theta, y) - \log q_\phi(\theta)\bigr]\dd\theta.
\]
Assuming we can exchange $\nabla_\phi\int = \int\nabla_\phi$, we get
\begin{align*}
\nabla_\phi\mathcal{L}(\phi) &= \int \nabla_\phi\Bigl(q_\phi(\theta)\bigl[\log p(\theta, y) - \log q_\phi(\theta)\bigr]\Bigr)\dd\theta \\
&= \int \bigl(\nabla_\phi q_\phi(\theta)\bigr)\bigl[\log p(\theta, y) - \log q_\phi(\theta)\bigr]\dd\theta
 + \int q_\phi(\theta)\,\nabla_\phi\bigl[\log p(\theta, y) - \log q_\phi(\theta)\bigr]\dd\theta.
\end{align*}
Since $\nabla_\phi\log p(\theta, y) = 0$,
\[
\nabla_\phi\mathcal{L}(\phi) = \int \bigl(\nabla_\phi q_\phi(\theta)\bigr)\bigl[\log p(\theta, y) - \log q_\phi(\theta)\bigr]\dd\theta
 - \underbrace{\int q_\phi(\theta)\,\nabla_\phi\log q_\phi(\theta)\dd\theta}_{=\,0}.
\]
The last term vanishes:
\[
\int q_\phi(\theta)\,\frac{1}{q_\phi(\theta)}\nabla_\phi q_\phi(\theta)\dd\theta
= \int \nabla_\phi q_\phi(\theta)\dd\theta
= \nabla_\phi\Bigl(\underbrace{\int q_\phi(\theta)\dd\theta}_{=\,1}\Bigr) = 0.
\]
Using the log-derivative trick,
\[
\frac{1}{q_\phi(\theta)}\nabla_\phi q_\phi(\theta) = \nabla_\phi\log q_\phi(\theta)
\quad\Longrightarrow\quad
\nabla_\phi q_\phi(\theta) = q_\phi(\theta)\,\nabla_\phi\log q_\phi(\theta),
\]
we obtain
\begin{align*}
\nabla_\phi\mathcal{L}(\phi) &= \int q_\phi(\theta)\bigl(\log p(\theta, y) - \log q_\phi(\theta)\bigr)\nabla_\phi\log q_\phi(\theta)\dd\theta \\
&= \E_{q_\phi(\theta)}\Bigl[\bigl(\log p(\theta, y) - \log q_\phi(\theta)\bigr)\nabla_\phi\log q_\phi(\theta)\Bigr].
\end{align*}

\subsection{Monte Carlo estimate}
With samples $\theta^{(s)} \sim q_\phi(\theta)$, $s = 1, \dots, S$,
\[
\widehat{\nabla_\phi\mathcal{L}}(\phi) = \frac{1}{S}\sum_{s=1}^{S}\Bigl[\log p\bigl(y, \theta^{(s)}\bigr) - \log q_\phi\bigl(\theta^{(s)}\bigr)\Bigr]\nabla_\phi\log q_\phi\bigl(\theta^{(s)}\bigr).
\]

%=====================================================
\section{Gaussian-mixture variational family}
For a mixture of Gaussians,
\[
q_\phi(\theta) = \sum_k \pi_k\,\Normal(\theta \given \mu_k, \Sigma_k),
\qquad
\log q_\phi(\theta) = \log\left[\sum_k \pi_k\,\Normal(\theta \given \mu_k, \Sigma_k)\right].
\]
Define
\[
x_k = \log\pi_k + \log\Normal(\theta \given \mu_k, \Sigma_k),
\]
so that
\[
e^{x_k} = e^{\log(\pi_k\,\Normal(\theta \mid \mu_k, \Sigma_k))} = \pi_k\,\Normal(\theta \given \mu_k, \Sigma_k),
\qquad
\sum_k e^{x_k} = \sum_k \pi_k\,\Normal(\theta \given \mu_k, \Sigma_k).
\]
Therefore
\[
\log q_\phi(\theta) = \log\sum_k e^{x_k},
\]
which is computed in a numerically stable way with the log-sum-exp function.

\end{document}
